\section{INTRODUCCIÓN}

\subsection{Antecedentes}

La vigilancia de patógenos virales en aguas residuales se ha consolidado como una estrategia complementaria para observar la circulación de agentes infecciosos a escala poblacional, pues, a diferencia de la vigilancia clínica, no depende del acceso a servicios de salud y de la notificación de casos. El análisis de aguas servidas integra señales biológicas de una comunidad conectada a una red de alcantarillado, un enfoque que se denomina epidemiología basada en aguas residuales. Esto permite generar información agregada sobre la presencia, la persistencia y los cambios temporales de microorganismos de interés sanitario \cite{Sims_2020}.

Durante la pandemia de COVID-19, la detección de ARN de SARS-CoV-2 en aguas residuales demostró que este tipo de muestra puede reflejar la circulación comunitaria del virus, incluso cuando la vigilancia clínica tiene limitaciones de cobertura o retrasos en la notificación \cite{Ahmed_2020,Medema_2020}. Además, la incorporación de métodos de secuenciación permitió pasar de la detección molecular hacia la vigilancia genómica, con capacidad para identificar linajes, variantes y diversidad viral presentes en una población \cite{Crits_Christoph_2021}. Este cambio amplió el valor de las aguas residuales como fuente de información para salud pública, especialmente en contextos urbanos donde múltiples sectores de la población convergen en una misma planta de tratamiento.

Además de SARS-CoV-2, las aguas residuales han sido utilizadas para la vigilancia ambiental de otros virus de interés sanitario, tales como poliovirus. Esta herramienta ha sido empleada dentro de los esfuerzos globales de erradicación para complementar la vigilancia de parálisis flácida aguda y detectar circulación silenciosa en comunidades \cite{Asghar_2014}. De forma similar, estudios recientes han mostrado que el ARN del virus sincitial respiratorio puede detectarse en aguas residuales y relacionarse con patrones o sospechas clínicas de infección respiratoria \cite{Hughes_2022}. Estos antecedentes sostienen la importancia de evaluar simultáneamente virus que puedan ser detectados en aguas servidas por la excreción de material viral por individuos infectados.

\subsection{Problemática}

La vigilancia de virus respiratorios y entéricos basada en casos clínicos depende de la presencia de síntomas, de la búsqueda de atención, del acceso al diagnóstico y de la notificación oportuna. Las revisiones sobre vigilancia en aguas residuales superan las limitaciones frecuentes de ese circuito: cobertura poblacional incompleta, disponibilidad limitada de pruebas, retrasos de reporte y falta de registro de infecciones leves o asintomáticas \cite{Bubba_2023,Parkins_2023}. En la pandemia de COVID-19, los casos notificados subestimaron la carga real debido a la limitada disponibilidad de pruebas para la detección del virus y a la reducción del análisis clínico durante la ola Ómicron BA.1, cuando la transmisión superó la capacidad diagnóstica disponible \cite{Parkins_2023}. Esta limitación disminuye la oportunidad de detectar la circulación comunitaria y reduce la priorización de las respuestas de salud pública.

La vigilancia basada en aguas residuales ofrece una visión representativa de la circulación viral en la población conectada al alcantarillado o a la planta de tratamiento. En Bangkok, en una investigación publicada en 2022, se analizaron 132 muestras de 19 plantas de tratamiento y se detectó ARN de SARS-CoV-2 aún cuando los casos reportados permanecían bajos; las cargas virales y los porcentajes de positividad se correlacionaron con los casos nuevos con rezagos de 22 a 24 días y valores de Spearman entre 0,85 y 1,00; lo cual indica una correlación positiva muy fuerte \cite{Sangsanont_2022}. En Calgary, investigadores estudiaron tres plantas de tratamiento y seis vecindarios durante un año, detectando SARS-CoV-2 en 406 de 414 muestras (98,06 \%) y correlacionándolo con casos clínicos a escala urbana, por planta y por barrio \cite{Acosta_2022}. En países de ingresos bajos y medios, estos programas enfrentan limitaciones en la infraestructura sanitaria, el financiamiento, la frecuencia de muestreo y la capacidad para la vigilancia de variantes \cite{Parkins_2023}.

La vigilancia clínica de poliovirus se apoya en la detección de parálisis flácida aguda, un desenlace que afecta a una fracción mínima de las personas infectadas: la tasa de ataque clínico del poliovirus virulento puede ser de un caso paralítico por cada 100 infecciones o menos \cite{Zurbriggen_2008}. En 2024, la detección de Poliovirus circulante derivado de la vacuna de tipo 2 (cVDPV2) en aguas residuales de Finlandia, Alemania, Polonia, España y Reino Unido mostró cepas genéticamente relacionadas y evidenció una circulación geográfica del virus más amplia de lo inicialmente detectado, lo que destacó el valor de la vigilancia ambiental para la detección temprana y el seguimiento de su dispersión \cite{Bottcher_2025}.

El virus sincitial respiratorio (RSV) plantea otro problema de vigilancia por su carga en lactantes, adultos mayores e inmunocomprometidos, y por su presentación clínica compartida con otros virus respiratorios. Hughes et al., 2022, detectaron ARN de RSV en sólidos sedimentados de aguas residuales de dos plantas y reportaron asociación fuerte con tasas de positividad clínica (Kendall tau = 0,65--0,77; $p < 10^{-7}$) \cite{Hughes_2022}. Le et al., 2025, desarrollaron un flujo de secuenciación ONT para muestras clínicas y aguas residuales, obtuvieron genomas completos en el 98 \% de las muestras con más de 10.000 copias/mL y observaron que las frecuencias de clados en aguas residuales representaban las frecuencias observadas en pacientes \cite{Le_2025}. El mismo estudio subraya la necesidad de una vigilancia genómica accesible en contextos con recursos limitados, debido a la menor representación de datos genómicos de RSV procedentes de países de ingresos bajos y medios \cite{Le_2025}.

En Quito, la PTAR Quitumbe ofrece un punto concreto para observar esta brecha. Mejía Calle (2024) realizó un muestreo pasivo semanal de 24 horas en el influente de la PTAR Quitumbe. En ese trabajo, los genes N1 y N2 de SARS-CoV-2 fueron detectados en todas las muestras, y la secuenciación con Oxford Nanopore permitió identificar linajes de Ómicron, incluida la detección de JN.1 en aguas residuales antes de su identificación clínica local \cite{MejiaCalle_2024}. Esta evidencia confirma la factibilidad del sitio para la vigilancia ambiental, que se centra en SARS-CoV-2.